\section*{Capítulo \ref{ch:b1:quantum-numbers} --- Números cuánticos y configuraciones electrónicas}

\begin{solution}{exo:b1:quantum-numbers:1}
(a) Imposible: $l \le n - 1 = 1$. (b) Posible (un electrón 3p).
(c) Imposible: $m_s = \pm\tfrac12$. (d) Posible (un electrón 4f).
(e) Imposible: para $l = 0$, $m_l$ solo puede valer 0.
\end{solution}

\begin{solution}{exo:b1:quantum-numbers:2}
3d: 5 orbitales, 10 electrones; 4f: 7 orbitales, 14 electrones; 5p: 3
orbitales, 6 electrones; capa $n = 2$: $n^2 = 4$ orbitales, $2n^2 = 8$
electrones.
\end{solution}

\begin{solution}{exo:b1:quantum-numbers:3}
\ce{O}: $1\text{s}^2 2\text{s}^2 2\text{p}^4 = [\ce{He}]\,2\text{s}^2 2\text{p}^4$,
6 \omterm{def:b1:quantum-numbers:configuration}{electrones de valencia}. \ce{P}: $1\text{s}^2 2\text{s}^2 2\text{p}^6 3\text{s}^2
3\text{p}^3 = [\ce{Ne}]\,3\text{s}^2 3\text{p}^3$, 5. \ce{K}: $[\ce{Ar}]\,4\text{s}^1$,
1. \ce{Ca}: $[\ce{Ar}]\,4\text{s}^2$, 2. \ce{Br}: $[\ce{Ar}]\,3\text{d}^{10}
4\text{s}^2 4\text{p}^5$, 7 (la \omterm{def:b1:quantum-numbers:orbital}{subcapa} 3d llena pertenece al núcleo).
% ledger: gs:K, gs:Ca
\end{solution}

\begin{solution}{exo:b1:quantum-numbers:4}
\ce{N}, 2p: \omorbs{u,u,u}, 3 desapareados. \ce{S}, 3p: \omorbs{ud,u,u},
2 desapareados. \ce{Cl}, 3p: \omorbs{ud,ud,u}, 1 desapareado. En los tres
casos, 2s o 3s está lleno, \omorbs{ud}.
\end{solution}

\begin{solution}{exo:b1:quantum-numbers:5}
\ce{Na+}, \ce{Mg^{2+}}, \ce{Al^{3+}}, \ce{F-}, \ce{O^{2-}}: $1\text{s}^2
2\text{s}^2 2\text{p}^6$, isoelectrónicos con el neón. \ce{Cl-}, \ce{S^{2-}},
\ce{Ca^{2+}}: $[\ce{Ne}]\,3\text{s}^2 3\text{p}^6$, isoelectrónicos con el
argón.
\end{solution}

\begin{solution}{exo:b1:quantum-numbers:6}
Los electrones 4s salen primero (\cref{prop:b1:quantum-numbers:ions}).
\ce{Ti^{2+}}: $[\ce{Ar}]\,3\text{d}^2$, 2 desapareados; \ce{Mn^{2+}}:
$3\text{d}^5$, 5; \ce{Co^{2+}}: $3\text{d}^7$, 3 (\omorbs{ud,ud,u,u,u});
\ce{Ni^{2+}}: $3\text{d}^8$, 2; \ce{Zn^{2+}}: $3\text{d}^{10}$, 0.
% ledger: gs:Ti2+, gs:Mn2+, gs:Co2+, gs:Zn2+
\end{solution}

\begin{solution}{exo:b1:quantum-numbers:7}
$\Delta E = 13.6\,(1/9 - 1/16) = 13.6 \times 7/144 = \qty{0.661}{eV}$;
$\lambda = 1240/0.661 = \qty{1876}{nm} \approx \qty{1.88}{\micro m}$:
infrarrojo (primera línea de la serie de Paschen).
\end{solution}

\begin{solution}{exo:b1:quantum-numbers:8}
(a) $[\ce{Ne}]\,3\text{s}^2 3\text{p}^4$: azufre, $Z = 16$.
(b) $[\ce{Ar}]\,3\text{d}^{10} 4\text{s}^2 4\text{p}^5$: bromo, $Z = 35$.
(c) $[\ce{Kr}]\,5\text{s}^1$: rubidio, $Z = 37$.
(d) $[\ce{Ar}]\,3\text{d}^3 4\text{s}^2$: vanadio, $Z = 23$.
Un electrón 3p del azufre: $(3, 1, -1, +\tfrac12)$, por ejemplo.
% ledger: gs:V
\end{solution}

\begin{solution}{exo:b1:quantum-numbers:9}
Para 4s, $n + l = 4 + 0 = 4$; para 3d, $n + l = 3 + 2 = 5$: 4s va
primero. Tras 3p ($n + l = 4$), la \omterm{def:b1:quantum-numbers:orbital}{subcapa} siguiente es 4s ($n + l = 4$,
$n$ mayor), que abre el periodo 4; 3d ($n + l = 5$) solo llega después
de 4s. El periodo 3 contiene, por tanto, solo 3s y 3p: $2 + 6 = 8$
elementos.
\end{solution}

\begin{solution}{exo:b1:quantum-numbers:10}
Energía de ionización $Z^2E_R = 4 \times 13.6 = \qty{54.4}{eV}$, cuatro
veces la del hidrógeno. $2 \to 1$: $\Delta E = 54.4\,(1 - 1/4) = \qty{40.8}{eV}$,
$\lambda = 1240/40.8 = \qty{30.4}{nm}$, cuatro veces más corta que la
línea del hidrógeno en \qty{121.6}{nm}: toda energía queda multiplicada
por $Z^2 = 4$.
\end{solution}

\begin{solution}{exo:b1:quantum-numbers:11}
(a) \omterm{def:b1:quantum-numbers:energy-level}{Estado excitado} del berilio (\omterm{def:b1:quantum-numbers:energy-level}{estado fundamental} $1\text{s}^2 2\text{s}^2$).
(b) Imposible: 2p admite como máximo 6 electrones. (c) \omterm{def:b1:quantum-numbers:energy-level}{Estado excitado} del
berilio. (d) \omterm{def:b1:quantum-numbers:energy-level}{Estado fundamental} del fósforo. (e) Imposible: una \omterm{def:b1:quantum-numbers:orbital}{subcapa} s
admite dos electrones (Pauli).
\end{solution}

\begin{solution}{exo:b1:quantum-numbers:12}
\ce{Fe^{3+}} $3\text{d}^5$: 5; \ce{Mn^{2+}} $3\text{d}^5$: 5; \ce{Fe^{2+}}
$3\text{d}^6$: 4; \ce{Cu^{2+}} $3\text{d}^9$: 1; \ce{Zn^{2+}}
$3\text{d}^{10}$: 0. Orden: $\ce{Fe^{3+}} = \ce{Mn^{2+}} > \ce{Fe^{2+}}
> \ce{Cu^{2+}} > \ce{Zn^{2+}}$. El hierro(III) tiene un \omterm{def:b1:quantum-numbers:unpaired}{electrón
desapareado} más que el hierro(II), de modo que es más atraído; el
cinc(II) no tiene ninguno, de modo que no es atraído.
% ledger: gs:Fe2+, gs:Fe3+, gs:Mn2+, gs:Cu2+, gs:Zn2+
\end{solution}

\begin{solution}{pb:b1:quantum-numbers:1}
\textbf{1.} $E_1 = \qty{-13.6}{eV}$, $E_2 = \qty{-3.40}{eV}$,
$E_3 = \qty{-1.51}{eV}$, $E_4 = \qty{-0.85}{eV}$.
\textbf{2.} $\Delta E = 3.40 - 1.51 = \qty{1.89}{eV}$, $\lambda = 1240/1.89
= \qty{656}{nm}$: rojo.
\textbf{3.} $4 \to 2$: \qty{2.55}{eV}, \qty{486}{nm} (verde azulado);
$5 \to 2$: \qty{2.86}{eV}, \qty{434}{nm} (violeta); $6 \to 2$: \qty{3.022}{eV},
\qty{410}{nm} (violeta).
\textbf{4.} $7 \to 2$: \qty{3.12}{eV}, \qty{397}{nm}, ya en el
ultravioleta; las transiciones superiores son aún más cortas, y $3 \to 2$
es la mayor longitud de onda de Balmer. Las líneas de Lyman están todas en
el ultravioleta y las de Paschen todas en el infrarrojo, de modo que solo
se ven cuatro líneas.
\textbf{5.} $\infty \to 2$: \qty{3.40}{eV}, $\lambda = \qty{365}{nm}$.
\textbf{6.} $\Delta E = 13.6 \times 3/4 = \qty{10.2}{eV}$, $\lambda =
\qty{122}{nm}$: ultravioleta lejano, invisible para el ojo y absorbido por
el vidrio y el aire.
\textbf{7.} El fotón debe aportar al menos \qty{13.6}{eV}: $\lambda \le
1240/13.6 = \qty{91.2}{nm}$.
\textbf{8.} $l = 0$ ($m_l = 0$), $l = 1$ ($m_l = -1, 0, 1$), $l = 2$
($m_l = -2, \dots, 2$): $1 + 3 + 5 = 9$ orbitales.
\textbf{9.} $\sum_{l=0}^{n-1}(2l+1) = n^2$ orbitales de dos electrones:
$2n^2$; para $n = 4$, 32.
\textbf{10.} 1s, 2s, 2p, 3s, 3p, 4s, 3d, 4p, 5s, 4d, 5p, 6s, 4f, 5d, 6p,
7s, 5f, 6d, 7p.
\textbf{11.} Cada periodo va de $n$s a $n$p: 2, 8, 8, 18, 18, 32, 32.
\textbf{12.} 3d ($n + l = 5$) va después de 4s ($n + l = 4$), que abre el
periodo 4.
\textbf{13.} $Z = 2, 10, 18, 36, 54, 86, 118$.
\textbf{14.} $[\ce{Ar}]\,3\text{d}^6 4\text{s}^2$; 3d: \omorbs{ud,u,u,u,u},
4 \omterm{def:b1:quantum-numbers:unpaired}{electrones desapareados}.
\textbf{15.} \ce{Fe^{2+}} $[\ce{Ar}]\,3\text{d}^6$: 4 desapareados;
\ce{Fe^{3+}} $[\ce{Ar}]\,3\text{d}^5$: 5 desapareados, una \omterm{def:b1:quantum-numbers:orbital}{subcapa} semillena.
\textbf{16.} Predicha $3\text{d}^4 4\text{s}^2$: 4 desapareados; observada
$3\text{d}^5 4\text{s}^1$: 6 desapareados (5 en 3d, 1 en 4s).
\textbf{17.} \ce{Cu} $[\ce{Ar}]\,3\text{d}^{10} 4\text{s}^1$, \ce{Cu+}
$[\ce{Ar}]\,3\text{d}^{10}$, \ce{Cu^{2+}} $[\ce{Ar}]\,3\text{d}^9$.
% ledger: gs:Cu, gs:Cu+, gs:Cu2+
\textbf{18.} \ce{Cu^{2+}}, con un \omterm{def:b1:quantum-numbers:unpaired}{electrón desapareado}; \ce{Cu+} no tiene
ninguno.
\textbf{19.} \ce{Sc^{3+}} $[\ce{Ar}]$: 0; \ce{Ti^{3+}} $[\ce{Ar}]\,3\text{d}^1$:
1; \ce{Mn^{2+}} $[\ce{Ar}]\,3\text{d}^5$: 5. \ce{Mn^{2+}} es el que más
tiene.
\textbf{20.} Tres electrones por orbital; $3(2l + 1)$ por \omterm{def:b1:quantum-numbers:orbital}{subcapa};
$3n^2$ por capa.
\textbf{21.} 1s: 3; 2s: 3; 2p: 9; 3s: 3; 3p: 9.
\textbf{22.} El primer periodo se cierra cuando 1s está lleno: $Z = 3$.
\textbf{23.} $Z = 4$, un electrón en 2s más allá del núcleo de gas noble.
\textbf{24.} 2s y 2p: $3 + 9 = 12$ elementos.
\textbf{25.} $3 + 12 = 15$: el segundo gas noble del mundo de tres espines
tiene \textbf{$Z = 15$}.
\end{solution}
